% Purpose: the measurements of the network: host and site contributions, the learning curve and the descriptor advantage (home of the 400-host, cross-validation-range and family-transfer limits), site resolution, staging and relaxed geometry; screening in sections/results_screening.tex. Mechanisms are left to the Discussion. Register: docs/prose_style.md. Every number is a macro of paper/generated/numbers.tex.
\section{Results}
\label{sec:results}

Errors are site-level mean absolute errors in eV on held-out oxides, intervals in square brackets are two-sided \pctci{} bootstrap intervals and $U$ is the one-sided upper bound of $A$ (Section~\ref{sec:statistics}); Table~\ref{tab:headline} collects the main values.

\subsection{Host and Site Contributions to the Formation Energy}
\label{sec:host_site}

Most of the variance of the vacancy energy lies between oxides. Grouped by host, the intraclass correlation of $E_\mathrm{f}$ \citep{Shrout1979intraclass} is $\mathrm{ICC}_\mathrm{host} = \numVarianceDecompositionHostIcc{}$ [\numVarianceDecompositionHostIccCiLo{}, \numVarianceDecompositionHostIccCiHi{}], with a between-host variance of \numVarianceDecompositionHostSigmaTwoBetween{}~\eVsq{} against \numVarianceDecompositionHostSigmaTwoWithin{}~\eVsq{} within hosts. We propose that the cost of removing a neutral oxygen atom is set mainly by the lattice it leaves, through the strength of its bonds and the band edges that receive the two electrons, and adjusted by the coordination of the site. The variance decomposition above shows that the variance lies between hosts; it does not establish which host property carries it, which Sections~\ref{sec:disc_principal} and~\ref{sec:disc_site} examine. A predictor that knows only the true mean of each host, the host-mean oracle, reaches \numHostMeanOracleAggregateMaeMean{} [\numHostMeanOracleAggregateMaeCiLo{}, \numHostMeanOracleAggregateMaeCiHi{}]~eV, and any predictor constant within a host has a within-host residual MAE of \numHostMeanOracleAggregateWithinHostMaeMean{}~eV, the zero-skill value.

Two host quantities already place oxides coarsely, and they saturate at once. The floor B0, linear in the oxide formation energy and the band gap, reaches \numNetTwentyFiveMaeBZero{}~eV with \numCfgBudgetsHostsItemOne{} training hosts and \numNetSixHundredFiftyFourMaeBZero{} [\numNetSixHundredFiftyFourMaeBZeroLo{}, \numNetSixHundredFiftyFourMaeBZeroHi{}]~eV with \numCfgBudgetsMax{}. RF-Kumagai reaches \numNetSixHundredFiftyFourMaeRf{} [\numNetSixHundredFiftyFourMaeRfLo{}, \numNetSixHundredFiftyFourMaeRfHi{}]~eV with \numCfgBudgetsMax{} training hosts, within \numBaselinesSummaryRfKumagaiVsPublishedDifferenceEV{}~eV of the published \numBaselinesSummaryRfKumagaiVsPublishedPublishedEV{}~eV \citep{Kumagai2021insights}. The distance between B0 and the host-mean oracle is the host-level variation they do not capture. B0 is constant within a host, so its within-host residual MAE equals the zero-skill value, \numBaselinesSummaryOuterFullBudgetBZeroWithinHostMaeMean{}~eV, whereas RF-Kumagai, which also reads the descriptors of the site, resolves part of the site-to-site variation, at \numBaselinesSummaryOuterFullBudgetRfKumagaiWithinHostMaeMean{} [\numBaselinesSummaryOuterFullBudgetRfKumagaiWithinHostMaeCiLo{}, \numBaselinesSummaryOuterFullBudgetRfKumagaiWithinHostMaeCiHi{}]~eV. At the same size the forest on the structural and compositional descriptors alone reaches \numBaselinesSummaryOuterFullBudgetRfStructuralMaeMean{}~eV and the forest on the electronic descriptors alone \numBaselinesSummaryOuterFullBudgetRfElectronicMaeMean{}~eV.

The structure-only network reproduces the published accuracy of structure-only prediction. With \numCfgBudgetsMax{} training hosts S reaches \numNetSixHundredFiftyFourMaeS{} [\numNetSixHundredFiftyFourMaeSLo{}, \numNetSixHundredFiftyFourMaeSHi{}]~eV on held-out oxides grouped by formula, against the \numNetPublishedMae{}~eV that Kiyohara et al.\ report on their own split \citep{Kiyohara2025machinelearning}.

\subsection{How Many Oxides Structure Alone Needs}
\label{sec:learning_curves}

The central question is at what number of training oxides the electronic structure of the pristine host stops lowering the site error by more than $\delta$. Both errors fall steeply as oxides are added (Figure~\ref{fig:learning_curves}a): S from \numNetTwentyFiveMaeS{}~eV with \numCfgBudgetsHostsItemOne{} training hosts to \numNetSixHundredFiftyFourMaeS{}~eV with \numCfgBudgetsMax{}, and D to \numNetSixHundredFiftyFourMaeD{}~eV, so the curves converge and $A$ shrinks (Figure~\ref{fig:learning_curves}b). At small sizes the descriptor forest is as accurate as either network or better. With \numCfgBudgetsHostsItemThree{} training hosts S lies above RF-Kumagai by \numNetCrossSMinusRfOneHundred{} [\numNetCrossSMinusRfOneHundredLo{}, \numNetCrossSMinusRfOneHundredHi{}]~eV and D matches it (\numNetCrossDMinusRfOneHundred{} [\numNetCrossDMinusRfOneHundredLo{}, \numNetCrossDMinusRfOneHundredHi{}]~eV). By \numCfgBudgetsHostsItemFour{} hosts D lies below the forest (\numNetCrossDMinusRfTwoHundred{} [\numNetCrossDMinusRfTwoHundredLo{}, \numNetCrossDMinusRfTwoHundredHi{}]~eV) and S below it in point estimate (\numNetCrossSMinusRfTwoHundred{} [\numNetCrossSMinusRfTwoHundredLo{}, \numNetCrossSMinusRfTwoHundredHi{}]~eV), with an interval below zero at \numNetCrossSBudget{} hosts (\numNetCrossSMinusRfSixHundredFiftyFour{} [\numNetCrossSMinusRfSixHundredFiftyFourLo{}, \numNetCrossSMinusRfSixHundredFiftyFourHi{}]~eV). With \numCfgBudgetsHostsItemOne{} and \numCfgBudgetsHostsItemTwo{} training hosts the error of D is dominated by fewer than one in a hundred test sites, in La and Ce oxides whose $f$ character of the conduction-band minimum lies far outside the training range (\numNetDsmallTwentyFiveExtremeShare{} and \numNetDsmallFiftyExtremeShare{} of the absolute error), so Figure~\ref{fig:learning_curves}b starts at \numCfgBudgetsHostsItemThree{} hosts.

\begin{figure}[H]
\centering
\figph{figure2_descriptor_advantage}
\caption{Site error and descriptor advantage against training hosts. (\textbf{a}) Test MAE of the structure-only network S and the descriptor network D on held-out oxides, with RF-Kumagai and the physics floor B0 as references, and (\textbf{b}) the advantage $A = \mathrm{MAE}_\mathrm{S} - \mathrm{MAE}_\mathrm{D}$ from \numCfgBudgetsHostsItemThree{} training hosts, with the tolerance $\pm\delta$, one-sided upper bounds, the additional resamples and grouped cross-validation at its mean training size. Bands, clipped in (\textbf{a}), and vertical bars are two-sided bootstrap intervals, values in \texttt{results/derived\_c1/final\_network\_summary.json}.\label{fig:learning_curves}}
\end{figure}

From \numNetNStar{} training oxides the structure-only network matches the descriptor network to within $\delta$. Over \numPrecisionABFourHundredBSixHundredFiftyFourSubsetsThirtyXFiveNResamples{} resamples with \numPrecisionABFourHundredBSixHundredFiftyFourSubsetsThirtyXFiveNSeeds{} seeds, $A = \numNetHeadlineA{}$ [\numNetHeadlineALo{}, \numNetHeadlineAHi{}]~eV with $U = \numNetHeadlineAUpper{}$~eV, and in the learning-curve design $A = \numNetSixHundredFiftyFourA{}$ [\numNetSixHundredFiftyFourALo{}, \numNetSixHundredFiftyFourAHi{}]~eV with $U = \numNetSixHundredFiftyFourAUpper{}$~eV; both intervals lie inside $\pm\delta$. With \numCfgBudgetsHostsItemFive{} training oxides the advantage is not bounded below $\delta$: $A = \numNetFourHundredA{}$ [\numNetFourHundredALo{}, \numNetFourHundredAHi{}]~eV with $U = \numNetFourHundredAUpper{}$~eV, and $U = \numNetFourHundredThirtyByFiveAUpper{}$~eV over the additional resamples, where \numNetSpreadFourHundredNAboveDelta{} of the \numCfgSplitNOuterResamples{} learning-curve resamples and \numNetSpreadFourHundredThirtyByFiveNAboveDelta{} of the \numPrecisionABFourHundredBSixHundredFiftyFourSubsetsThirtyXFiveNResamples{} resamples have $A$ above $\delta$. With \numCfgBudgetsHostsItemFour{} training hosts the bound falls to \numNetTwoHundredAUpper{}~eV because the point estimate there is low, \numNetTwoHundredA{}~eV, at the same precision (bootstrap standard deviation \numNetSpreadTwoHundredBootSd{}~eV, against \numNetSpreadFourHundredBootSd{}~eV with \numCfgBudgetsHostsItemFive{}). $N^{*}$, the smallest size from which the bound holds at every larger size, is therefore \numNetNStar{}.

The match holds beyond the learning curve. Grouped cross-validation, training on about \numNetCvTenHosts{} and \numNetCvTwentyHosts{} hosts per fold, gives $A = \numNetCvTenA{}$ [\numNetCvTenALo{}, \numNetCvTenAHi{}] and \numNetCvTwentyA{} [\numNetCvTwentyALo{}, \numNetCvTwentyAHi{}]~eV with $U = \numNetCvTenAUpper{}$ and \numNetCvTwentyAUpper{}~eV, both intervals inside $\pm\delta$. Cross-validation tests every host once, including hosts that no resample of the learning curve tests, so its points describe a different test population, and nothing is known beyond about \numNetCvTwentyHosts{} training hosts. From \numNetNStar{} training oxides on, the advantage stays positive, with intervals that exclude zero in the additional resamples and in both cross-validations: the electronic structure of the pristine host still lowers the site error, by at most \numNetHeadlineAUpper{}~eV.

\begin{table}[H]
\caption{Main results of the network. Site-level MAE or advantage $A = \mathrm{MAE}_\mathrm{S} - \mathrm{MAE}_\mathrm{D}$ in eV on held-out oxides, with the two-sided \pctci{} interval in brackets and, for $A$, the one-sided upper bound $U$; $\delta = \numNetTolerance{}$~eV. $B$ is the number of training hosts. Values in \texttt{results/derived\_c1/final\_network\_summary.json} and the records of Appendix~\ref{app:code}.\label{tab:headline}}
\centering
\small
\begin{tabularx}{\linewidth}{Xl}
\toprule
Quantity & Value (eV) \\
\midrule
$\mathrm{ICC}_\mathrm{host}$ of $E_\mathrm{f}$ (dimensionless) & \numVarianceDecompositionHostIcc{} [\numVarianceDecompositionHostIccCiLo{}, \numVarianceDecompositionHostIccCiHi{}] \\
Host-mean oracle, MAE & \numHostMeanOracleAggregateMaeMean{} [\numHostMeanOracleAggregateMaeCiLo{}, \numHostMeanOracleAggregateMaeCiHi{}] \\
B0, MAE, $B = \numCfgBudgetsMax{}$ & \numNetSixHundredFiftyFourMaeBZero{} [\numNetSixHundredFiftyFourMaeBZeroLo{}, \numNetSixHundredFiftyFourMaeBZeroHi{}] \\
RF-Kumagai, MAE, $B = \numCfgBudgetsMax{}$ & \numNetSixHundredFiftyFourMaeRf{} [\numNetSixHundredFiftyFourMaeRfLo{}, \numNetSixHundredFiftyFourMaeRfHi{}] \\
S, MAE, $B = \numCfgBudgetsMax{}$ & \numNetSixHundredFiftyFourMaeS{} [\numNetSixHundredFiftyFourMaeSLo{}, \numNetSixHundredFiftyFourMaeSHi{}] \\
$A$, $B = \numCfgBudgetsHostsItemFive{}$, learning curve & \numNetFourHundredA{} [\numNetFourHundredALo{}, \numNetFourHundredAHi{}]; $U$ \numNetFourHundredAUpper{} \\
$A$, $B = \numCfgBudgetsHostsItemFive{}$, additional resamples & \numNetFourHundredThirtyByFiveA{} [\numNetFourHundredThirtyByFiveALo{}, \numNetFourHundredThirtyByFiveAHi{}]; $U$ \numNetFourHundredThirtyByFiveAUpper{} \\
$A$, $B = \numCfgBudgetsMax{}$, learning curve & \numNetSixHundredFiftyFourA{} [\numNetSixHundredFiftyFourALo{}, \numNetSixHundredFiftyFourAHi{}]; $U$ \numNetSixHundredFiftyFourAUpper{} \\
$A$, $B = \numCfgBudgetsMax{}$, additional resamples & \numNetHeadlineA{} [\numNetHeadlineALo{}, \numNetHeadlineAHi{}]; $U$ \numNetHeadlineAUpper{} \\
$A$, cross-validation, $K = \numSplitCvKsItemOne{}$ & \numNetCvTenA{} [\numNetCvTenALo{}, \numNetCvTenAHi{}]; $U$ \numNetCvTenAUpper{} \\
$A$, cross-validation, $K = \numSplitCvKsItemTwo{}$ & \numNetCvTwentyA{} [\numNetCvTwentyALo{}, \numNetCvTwentyAHi{}]; $U$ \numNetCvTwentyAUpper{} \\
$A$, leave one family out & \numNetLocoA{} [\numNetLocoALo{}, \numNetLocoAHi{}]; $U$ \numNetLocoAUpper{} \\
$\mathrm{MAE}_\mathrm{P} - \mathrm{MAE}_\mathrm{S}$, $B = \numCfgBudgetsMax{}$, at the tuning and training length used & \numNetStagingPMinusSSixHundredFiftyFour{} [\numNetStagingPMinusSSixHundredFiftyFourLo{}, \numNetStagingPMinusSSixHundredFiftyFourHi{}] \\
$\Delta\mathrm{MAE}$ of S, relaxed geometry, $B = \numCfgBudgetsMax{}$ & \numNetGeomDeltaMaeSSixHundredFiftyFourRelaxed{} [\numNetGeomDeltaMaeSSixHundredFiftyFourRelaxedLo{}, \numNetGeomDeltaMaeSSixHundredFiftyFourRelaxedHi{}] \\
\bottomrule
\end{tabularx}
\end{table}

Transfer to cation chemistry absent from training is the exception. When whole families of cation chemistry are held out, $A = \numNetLocoA{}$ [\numNetLocoALo{}, \numNetLocoAHi{}]~eV with $U = \numNetLocoAUpper{}$~eV, above $\delta$. One family carries the excess, the alkali tetrelates (silicates, germanates, stannates and plumbates of the alkali metals; \numNetLocoAlkaliGroupFourteenHosts{} hosts), within which S errs by \numNetLocoAlkaliGroupFourteenAFamily{}~eV more than D per site; without that family the bound is \numNetLocoUpperWithoutAlkaliGroupFourteen{}~eV.

\subsection{Site Resolution within a Host}
\label{sec:within_host}

The local chemistry of the vacancy shows only in how well a model tells the oxygen sites of one oxide apart. Both networks partly resolve sites with \numCfgBudgetsMax{} training hosts: the within-host residual MAE is \numCOneCTwoVTwoMaeSixHundredFiftyFourSWithinHostMaeMean{}~eV for S and \numCOneCTwoVTwoMaeSixHundredFiftyFourDStateWithinHostMaeMean{}~eV for D, below the zero-skill value of \numHostMeanOracleAggregateWithinHostMaeMean{}~eV, a within-host skill of \numCOneCTwoVTwoMaeSixHundredFiftyFourSWithinHostSkillMean{} [\numCOneCTwoVTwoMaeSixHundredFiftyFourSWithinHostSkillCiLo{}, \numCOneCTwoVTwoMaeSixHundredFiftyFourSWithinHostSkillCiHi{}] and \numCOneCTwoVTwoMaeSixHundredFiftyFourDStateWithinHostSkillMean{} [\numCOneCTwoVTwoMaeSixHundredFiftyFourDStateWithinHostSkillCiLo{}, \numCOneCTwoVTwoMaeSixHundredFiftyFourDStateWithinHostSkillCiHi{}]~eV (Figure~\ref{fig:within_host}a). D spreads its predictions within a host more widely than S, to \numNetSpreadRatioDSixHundredFiftyFour{} of the DFT spread against \numNetSpreadRatioSSixHundredFiftyFour{}, without placing the sites more accurately. Within a host the descriptors give D no advantage (Figure~\ref{fig:within_host}b): $A_\mathrm{within} = \numNetWithinA{}$ [\numNetWithinALo{}, \numNetWithinAHi{}]~eV with \numCfgBudgetsMax{} training hosts, its upper bound lies below $\delta$ at every size ($N^{*}_\mathrm{within} = \numNetWithinNStar{}$), and the two cross-validations give \numCvExtendedBudgetVTwoKTenAWithinMean{} and \numCvExtendedBudgetVTwoKTwentyAWithinMean{}~eV. Over the \numNetWithinSpearmanHosts{} hosts with at least three tested sites, S and D rank the sites alike (difference of the mean within-host Spearman correlations \numNetWithinSpearmanSMinusD{} [\numNetWithinSpearmanSMinusDLo{}, \numNetWithinSpearmanSMinusDHi{}]), and the Bader charge of the oxygen to be removed ranks them with a mean correlation of only \numWithinHostBreakdownSpearmanReferenceBaderChargeMean{} [\numWithinHostBreakdownSpearmanReferenceBaderChargeCiNinetyFiveItemOne{}, \numWithinHostBreakdownSpearmanReferenceBaderChargeCiNinetyFiveItemTwo{}]. The site-to-site spread itself is larger in hosts with a filled cation $d$ shell: the zero-skill value, averaged over sites, is \numWithinHostBreakdownGroupsDClassDZeroVTwoSZeroSkillMaePoint{}~eV in hosts whose cations have an empty $d$ shell and \numWithinHostBreakdownGroupsDClassDTenVTwoSZeroSkillMaePoint{}~eV in hosts with a filled $d$ shell (Appendix~\ref{app:within}).

\begin{figure}[H]
\centering
\figph{figure3_within_host}
\caption{Site resolution within a host against training hosts. (\textbf{a}) Within-host skill of S and D, the zero-skill within-host residual MAE minus that of the network, on held-out hosts with at least two oxygen sites, and (\textbf{b}) the within-host advantage $A_\mathrm{within}$ with the tolerance $\pm\delta$. Bands and vertical bars are two-sided bootstrap intervals for the learning curve, the additional resamples and grouped cross-validation.\label{fig:within_host}}
\end{figure}

\subsection{Predicting the Descriptors from Structure}
\label{sec:staging}

If the descriptors could be predicted from structure, their value below \numNetNStar{} oxides would come without the DFT step. The staged model P tests this through the identity $\mathrm{MAE}_\mathrm{P} - \mathrm{MAE}_\mathrm{S} = -A - (\mathrm{MAE}_\mathrm{D} - \mathrm{MAE}_\mathrm{P})$: staging beats S only where the descriptor advantage exceeds the cost of predicting the descriptors. At the tuning and training length used, P is never better than S (Figure~\ref{fig:staging}); P1 kept the settings tuned under the default initialisation (Section~\ref{sec:training}). The outcome is inconclusive with \numCfgBudgetsHostsItemOne{} training hosts, P worse from \numCfgBudgetsHostsItemTwo{} to \numCfgBudgetsHostsItemFour{} (\numNetStagingPMinusSTwoHundred{} [\numNetStagingPMinusSTwoHundredLo{}, \numNetStagingPMinusSTwoHundredHi{}]~eV with \numCfgBudgetsHostsItemFour{}), equivalent with \numCfgBudgetsHostsItemFive{} (\numNetStagingPMinusSFourHundred{} [\numNetStagingPMinusSFourHundredLo{}, \numNetStagingPMinusSFourHundredHi{}]~eV), and P worse within the margin $\delta_\mathrm{P}$ with \numCfgBudgetsMax{}, where $\mathrm{MAE}_\mathrm{P} - \mathrm{MAE}_\mathrm{S} = \numNetStagingPMinusSSixHundredFiftyFour{}$ [\numNetStagingPMinusSSixHundredFiftyFourLo{}, \numNetStagingPMinusSSixHundredFiftyFourHi{}]~eV and the cost of prediction is $\mathrm{MAE}_\mathrm{D} - \mathrm{MAE}_\mathrm{P} = \numNetStagingDMinusPSixHundredFiftyFour{}$ [\numNetStagingDMinusPSixHundredFiftyFourLo{}, \numNetStagingDMinusPSixHundredFiftyFourHi{}]~eV. With \numCfgBudgetsHostsItemOne{} and \numCfgBudgetsHostsItemTwo{} training hosts D is itself unstable (Section~\ref{sec:learning_curves}), so the cost of staging is not interpretable there.

\begin{figure}[H]
\centering
\figph{figure4_staging}
\caption{Staged prediction through descriptors against training hosts. (\textbf{a}) Test MAE of S, D and the staged model P, which passes the descriptors predicted by P1 from the pristine structure to D of the same training set and seed, and (\textbf{b}) the cost of prediction $\mathrm{MAE}_\mathrm{D} - \mathrm{MAE}_\mathrm{P}$ and the staging contrast $\mathrm{MAE}_\mathrm{P} - \mathrm{MAE}_\mathrm{S}$ with the margins $\pm\delta_\mathrm{P}$. Bands are two-sided bootstrap intervals, clipped at the axis limits.\label{fig:staging}}
\end{figure}

The descriptor predictor recovers host-level electronic structure better than site-level. With \numCfgBudgetsMax{} training hosts the coefficient of determination of P1 on held-out hosts, averaged over each class, is \numStagingVTwoPOneDescriptorsSixHundredFiftyFourByClassHostElectronicClassMeanRTwo{} for the host-electronic and \numStagingVTwoPOneDescriptorsSixHundredFiftyFourByClassSiteElectronicClassMeanRTwo{} for the site-electronic descriptors; the oxide formation energy is recovered best, at \numStagingVTwoPOneDescriptorsSixHundredFiftyFourPerDescriptorDescFormationEnergyRTwoMean{}, and the Bader charge of the oxygen to be removed reaches only \numStagingVTwoPOneDescriptorsSixHundredFiftyFourPerDescriptorDescBaderChargeRTwoMean{}. P1 reaches the training-length limit of \numNetCapMaxEpochs{} epochs in \numNetCapPOneTwoHundred{}, \numNetCapPOneFourHundred{} and \numNetCapPOneSixHundredFiftyFour{} of \numNetCapRunsSixHundredFiftyFour{} runs with \numCfgBudgetsHostsItemFour{}, \numCfgBudgetsHostsItemFive{} and \numCfgBudgetsMax{} training hosts (Appendix~\ref{app:sens}).

\subsection{Relaxed Host Geometry}
\label{sec:geometry}

A structure relaxed with a machine-learned interatomic potential barely changes the predictions. MACE-MP-0 changes each host in three ways: it expands the cell (median isotropic strain \numNetGeomIsotropicStrainMedian{}), changes its shape at fixed volume (median deviatoric strain \numNetGeomDeviatoricStrainMedian{}) and displaces the atoms within it, which alters M--O bond lengths and polyhedral distortions (median \numNetGeomInternalRmsdMedian{}~\AA{}), over the \numNetGeomEvaluatedHosts{} hosts that were relaxed and evaluated. With \numCfgBudgetsMax{} training hosts the error of S changes by \numNetGeomDeltaMaeSSixHundredFiftyFourRelaxed{} [\numNetGeomDeltaMaeSSixHundredFiftyFourRelaxedLo{}, \numNetGeomDeltaMaeSSixHundredFiftyFourRelaxedHi{}]~eV on the relaxed cells and by \numNetGeomDeltaMaeSSixHundredFiftyFourRescaled{} [\numNetGeomDeltaMaeSSixHundredFiftyFourRescaledLo{}, \numNetGeomDeltaMaeSSixHundredFiftyFourRescaledHi{}]~eV on the cells rescaled to the DFT volume. The error of D changes by \numNetGeomDeltaMaeDSixHundredFiftyFourRelaxed{} [\numNetGeomDeltaMaeDSixHundredFiftyFourRelaxedLo{}, \numNetGeomDeltaMaeDSixHundredFiftyFourRelaxedHi{}]~eV and that of P by \numNetGeomDeltaMaePSixHundredFiftyFourRelaxed{} [\numNetGeomDeltaMaePSixHundredFiftyFourRelaxedLo{}, \numNetGeomDeltaMaePSixHundredFiftyFourRelaxedHi{}]~eV on the relaxed cells (Appendix~\ref{app:geometry}).

% Purpose: Results 3.6, the screening of named oxides with the network: absent oxides as a domain result, the error over the screened sites and its misses, the two energy groups and the within-host pairs. Input by sections/results.tex. Every number is a macro (records results/derived_w2/w2_screening_sites.json, figures/main/data/screening_demo.json, results/screening_absence.json).
\subsection{Screening of Named Oxides}
\label{sec:screening}

The last question is how close the structure-only network comes to DFT for oxides that a defect chemist would name, and to which of them it applies at all. This is a demonstration on named oxides: S with \numNetScreenBudget{} training hosts is applied to the \numScreeningDemoScopeNCandidates{} candidate oxides of the analysis plan, and the prediction for a site is the seed-ensemble mean over the resamples that hold its host out.

The absences are a result about the domain. Only \numScreeningDemoScopeNPresent{} of the \numScreeningDemoScopeNCandidates{} candidates are in the dataset, with \numNetScreenSites{} oxygen sites in \numNetScreenHosts{} hosts (Table~\ref{tab:screening_sites}). In$_2$O$_3$, CeO$_2$, WO$_3$, MoO$_3$, Nb$_2$O$_5$ and LaAlO$_3$ have no host of that formula in the release. SrTiO$_3$ and BaTiO$_3$ are in the release, but every neutral entry carries the perturbed-host-state flag (\numScreeningDemoScopePhsRemovedSrTiOThreePhsEntries{} of \numScreeningDemoScopePhsRemovedSrTiOThreeNeutralEntries{} and \numScreeningDemoScopePhsRemovedBaTiOThreePhsEntries{} of \numScreeningDemoScopePhsRemovedBaTiOThreeNeutralEntries{}), so the filter of Section~\ref{sec:data} removes them, and the monoclinic TiO$_2$ host keeps one oxygen site because \numScreeningDemoScopeTioTwoReleasePhsEntries{} of its \numScreeningDemoScopeTioTwoReleaseNeutralEntries{} neutral entries carry the flag. Every candidate that the filter removes, in whole or in part, is a Ti oxide, whose empty $d$ band lies low.

\begin{table}[H]
\caption{Screened oxygen sites. Cations: number of cations in the first coordination shell of the oxygen to be removed. DFT: formation energy of the release. S: prediction of the structure-only network with \numNetScreenBudget{} training hosts, the mean over the $N$ resamples in which the host was held out. Values in \texttt{results/derived\_w2/w2\_screening\_sites.json}, full columns in Appendix~\ref{app:screening}.\label{tab:screening_sites}}
\centering
\small
\begin{tabular}{llrrrr}
\toprule
Oxide & Site & Cations & DFT (eV) & S (eV) & $N$ \\
\midrule
SnO$_2$ & O$_1$ & \numNetScreenSnOTwoOOneCations{} & \numNetScreenSnOTwoOOneDft{} & \numNetScreenSnOTwoOOnePred{} & \numNetScreenSnOTwoOOneN{} \\
ZnO & O$_1$ & \numNetScreenZnOOOneCations{} & \numNetScreenZnOOOneDft{} & \numNetScreenZnOOOnePred{} & \numNetScreenZnOOOneN{} \\
Ga$_2$O$_3$ & O$_1$ & \numNetScreenGaTwoOThreeOOneCations{} & \numNetScreenGaTwoOThreeOOneDft{} & \numNetScreenGaTwoOThreeOOnePred{} & \numNetScreenGaTwoOThreeOOneN{} \\
Ga$_2$O$_3$ & O$_2$ & \numNetScreenGaTwoOThreeOTwoCations{} & \numNetScreenGaTwoOThreeOTwoDft{} & \numNetScreenGaTwoOThreeOTwoPred{} & \numNetScreenGaTwoOThreeOTwoN{} \\
Ga$_2$O$_3$ & O$_3$ & \numNetScreenGaTwoOThreeOThreeCations{} & \numNetScreenGaTwoOThreeOThreeDft{} & \numNetScreenGaTwoOThreeOThreePred{} & \numNetScreenGaTwoOThreeOThreeN{} \\
TiO$_2$ & O$_4$ & \numNetScreenTiOTwoOFourCations{} & \numNetScreenTiOTwoOFourDft{} & \numNetScreenTiOTwoOFourPred{} & \numNetScreenTiOTwoOFourN{} \\
ZrO$_2$ & O$_1$ & \numNetScreenZrOTwoOOneCations{} & \numNetScreenZrOTwoOOneDft{} & \numNetScreenZrOTwoOOnePred{} & \numNetScreenZrOTwoOOneN{} \\
ZrO$_2$ & O$_2$ & \numNetScreenZrOTwoOTwoCations{} & \numNetScreenZrOTwoOTwoDft{} & \numNetScreenZrOTwoOTwoPred{} & \numNetScreenZrOTwoOTwoN{} \\
MgO & O$_1$ & \numNetScreenMgOOOneCations{} & \numNetScreenMgOOOneDft{} & \numNetScreenMgOOOnePred{} & \numNetScreenMgOOOneN{} \\
Al$_2$O$_3$ & O$_1$ & \numNetScreenAlTwoOThreeOOneCations{} & \numNetScreenAlTwoOThreeOOneDft{} & \numNetScreenAlTwoOThreeOOnePred{} & \numNetScreenAlTwoOThreeOOneN{} \\
\bottomrule
\end{tabular}
\end{table}

Over the \numNetScreenSites{} screened sites the MAE is \numNetScreenMae{} [\numNetScreenMaeLo{}, \numNetScreenMaeHi{}]~eV, with the interval from a bootstrap over the \numNetScreenHosts{} hosts, against \numNetSixHundredFiftyFourMaeS{}~eV over all held-out sites; these oxides are not a random sample. The errors carry no net sign (bias \numNetScreenBias{}~eV; \numNetScreenOver{} sites over-predicted and \numNetScreenUnder{} under-predicted). The largest miss is the O$_2$ site of ZrO$_2$, predicted \numNetScreenZrOTwoOTwoErr{}~eV too high, and MgO and Al$_2$O$_3$ are predicted \numNetScreenMgOOOneAbsErr{} and \numNetScreenAlTwoOThreeOOneAbsErr{}~eV too low.

\begin{figure}[H]
\centering
\figph{figure5_screening}
\caption{Formation energies of the screened oxygen sites predicted by S with \numNetScreenBudget{} training hosts. (\textbf{a}) Prediction against the DFT energy, coloured by oxide, and (\textbf{b}) predicted minus DFT energy by oxide, in ascending order of the mean DFT energy of the oxide. Markers are means over the resamples that hold the host out and bars run from the smallest to the largest prediction, values in \texttt{figures/main/data/screening\_demo.json}.\label{fig:screening}}
\end{figure}

Across oxides the network reproduces the energy scale. The DFT energies fall into two groups (Figure~\ref{fig:screening}a): the sites of SnO$_2$, ZnO, Ga$_2$O$_3$ and TiO$_2$ lie between \numNetScreenGaTwoOThreeOOneDft{} and \numNetScreenTiOTwoOFourDft{}~eV, and those of ZrO$_2$, MgO and Al$_2$O$_3$ between \numNetScreenZrOTwoOTwoDft{} and \numNetScreenAlTwoOThreeOOneDft{}~eV. The network places every site in its group, and the rank correlation of predicted with DFT energy, \numNetScreenSpearman{}, comes from that separation. A structure render of the SnO$_2$ site is given in Appendix~\ref{app:screening}.

Within a host, the network orders \numNetScreenPairsConcordant{} of the \numNetScreenPairs{} site pairs as DFT does. In Ga$_2$O$_3$ the four-coordinate site O$_2$ is the highest, as in DFT, but the difference between the two three-coordinate sites, \numScreeningDemoPairsGaTwoOThreeOOneOThreeDftDiffEV{}~eV in DFT, shrinks to \numScreeningDemoPairsGaTwoOThreeOOneOThreeVTwoDiffEV{}~eV. In ZrO$_2$ the network inverts the order of the three-coordinate site O$_1$ and the four-coordinate site O$_2$, whose DFT difference is \numScreeningDemoPairsZrOTwoOOneOTwoDftDiffEV{}~eV and predicted difference \numScreeningDemoPairsZrOTwoOOneOTwoVTwoDiffEV{}~eV.

